\section{Appendix: Scope and Qualifications}

The First Sunday of Advent in the universal 1962 Roman calendar falls on 29 November in 2026. It is a first-class Sunday in violet with the Roman station at St Mary Major. The Creed is said, the Gloria in excelsis omitted, and the Preface of the Most Holy Trinity prescribed. St Saturninus's ordinary commemoration on the same date is omitted under the general rubrics, §§109 and 111(a). The Roman Canon has its ordinary form. The weekday reuse of the Sunday Mass retains the Gradual alone, omits the Creed and has the common Preface. No local calendar or institutional addition is presumed here.\footnote{1962 Missal, \work{Rubricae generales} §§71, 109, 111(a), and \emph{Proprium de tempore}, pp. [1]--3, including the continuation of the ferial rubric. For St Saturninus, \emph{Proprium sanctorum}, p. [428], nos. 1809--1811. The FSSP Province de France and ICRSP Province de France Ordos, entries for 29 November 2026, corroborate the Sunday. The ICRSP's papal prayer belongs to its own notice.}

The bilingual texts comprise the eleven appointments for this Mass. The 1962 typical edition controls the Latin text and the rubrical choices. The inherited Latin was compared with Pustet's 1862 missal. The Lasance and Walsh English belongs to the revised 1945 edition in its later printing. The U.S. public-domain basis for the English applies to this material outside that printing's later inserts. The two historical missals preserve earlier discipline concerning additional prayers. The modern page images, digital wrappers and editorial matter are outside the grant for the inherited wording. The historical texts have their own rights status, separate from the study's license.

Challoner's biblical English follows the registered Project Gutenberg edition. In the scriptural appointments it is given to the extent used at Mass, with historical liturgical frames distinguished from biblical words. The chants have their own adaptations. The two English witnesses preserve a word for time in Rom 13:11 where the Latin lesson still has the hour but omits the preceding biblical noun. At the Gradual, Shepherd's English preserves the added vocative. At the Communion the initial biblical connective is omitted in keeping with the chant. The Lasance English prayer endings printed on the Sunday are cues. The full forms and Amens printed alongside them are from the hand missal's table of conclusions. These historical study texts are aids to reading the Latin.

Augustine, Gregory, Chrysostom, Ambrose, Bede, Thomas and Bellarmine are cited for explanations of particular passages or doctrines. The comparisons extending those explanations across this Mass are the author's. Guéranger comments directly on the Sunday in his own books. His additional prayers do not apply to the universal-calendar Mass studied here. His traditional attribution of the Communion verse to David does not displace the superscription naming the sons of Core. The two senses of the fruit in Ambrose and Bede's two possibilities for the generation remain alternatives. Augustine's applications of the Psalms turn at points on wording different from the Missal's. The chapters and verses under consideration have their own evidence of place and date in \hyperref[sec:date-location]{Scriptural Date and Location}.

\section{References}

\begin{itemize}
\item \textbf{Missal.} \work{Missale Romanum}, Vatican typical edition, 1962, First Sunday of Advent, pp. [1]--3, nos. 1--10, and Trinity Preface, p. 293, no. 1082. \url{https://media.churchmusicassociation.org/pdf/missale62.pdf}. Latin antecedent in \work{Missale Romanum} (Ratisbon, Pustet, 1862), pp. 1--3, 251--252. \url{https://archive.org/details/bub_gb_E7sPAAAAIAAJ}. Francis Xavier Lasance and Francis Augustine Walsh, \work{The New Roman Missal} (Benziger, revised 1945), First Sunday of Advent, pp. 103--106, with the dialogue, Preface and Sanctus, pp. 772--774, and Prayer Conclusions, p. 819. \url{https://archive.org/details/fr-lasance-missal-1937}.
\item \textbf{Scripture.} Douay--Rheims, Challoner revision, per-book verse texts in the Project Gutenberg composite electronic edition, for the English of the biblical appointments and their contexts. \url{https://www.gutenberg.org/ebooks/1581}. Clementine Vulgate, eBible \emph{latVUC}, for comparison of the biblical text with the Missal.
\item \textbf{Calendar.} Priestly Fraternity of St Peter, Province de France, \work{Ordo du mois}, November 2026, and Institute of Christ the King Sovereign Priest, Province de France, \work{Ordo liturgique}, 29 November 2026. \url{https://www.fssp.fr/ordo-du-mois/} and \url{https://icrspfrance.fr/ordo.php?d=2026-11-29}.
\item \textbf{Augustine of Hippo.} \work{Enarrationes in Psalmos}, historical English in \work{Nicene and Post-Nicene Fathers}, first series, vol. VIII. English headings Psalm XXV, XLVIII, LXXX and LXXXV correspond to the Missal's Psalms 24, 47, 79 and 84. Section numbers in the notes are those of the historical English. \work{De Trinitate} III, preface §3, trans. Arthur West Haddan, on the missions of the Son and Spirit.
\item \textbf{John Chrysostom.} \work{Homilies on Romans}, Homily 24, commentary on the end of Rom 13. Historical English in \work{Nicene and Post-Nicene Fathers}, first series, vol. XI, Romans translation credited to J. B. Morris and W. H. Simcox. The New Advent edition also gives the text under Homily XXIV. \url{https://www.newadvent.org/fathers/210224.htm}.
\item \textbf{Gregory the Great.} \work{Homiliae in Evangelia}, book I, Homily 1, §§1--5, Migne-derived Latin text of the homilies. The transmitted pericope heading ends at verse 32; the text and exposition include Lk 21:33.
\item \textbf{Ambrose of Milan.} \work{Expositio evangelii secundum Lucam}, X.39--45, Latin text, on the powers, clouds and trees in Luke's discourse.
\item \textbf{Bede.} \work{In Lucae evangelium expositio}, book VI, at Lk 21:32, in Migne, \work{Patrologia Latina} 92, col. 590B--C.
\item \textbf{Thomas Aquinas.} \work{Super Epistolam B. Pauli ad Romanos lectura}, cap. 13, lectio 3, Corpus Thomisticum [86224], Enrique Alarcón's revision of Roberto Busa's Latin text from the 1953 Turin printing. \url{https://www.corpusthomisticum.org/cro12.html}.
\item \textbf{Robert Bellarmine.} \work{Commentary on the Book of Psalms}, John O'Sullivan's abridged English (Dublin, James Duffy, 1866), Ps 24 on the opening petition and Ps 84, exposition 12, p. 259, on the promise of fruit. The paragraph number in this edition is one less than the Douay verse number for the Communion.
\item \textbf{Prosper Guéranger.} \work{The Liturgical Year}, vol. I, Advent, trans. Laurence Shepherd (Dublin, James Duffy, 1870), First Sunday of Advent, especially pp. 134--137. \url{https://archive.org/details/liturgicalyear01gura}.
\end{itemize}
